\documentclass[12pt]{article}
\usepackage[utf8]{inputenc}
\usepackage[T1]{fontenc}
\usepackage[margin=1in]{geometry}
\usepackage{setspace}
\usepackage{xcolor}
\usepackage[colorlinks=true, urlcolor=blue!60!black, linkcolor=blue!60!black]{hyperref}
\setlength{\parindent}{2em}
\setlength{\parskip}{1.25em}
\renewcommand{\baselinestretch}{1.0}
\begin{document}
\begin{center}
{\LARGE \textit{Central Hudson Gas \& Electric Corp. v. Public Service Commission}, 447 U.S. 557 (1980)\par}
\vspace{0.8em}
{\large Case Brief --- Four-Page Edition\par}
\end{center}
\vspace{1em}

\section{Memory Jogger}

Striking down New York's total ban on promotional advertising by electric utilities 8--1, the Court announced a four-part test: whether the speech concerns lawful activity and is not misleading, whether the government asserts a substantial interest, whether the restriction directly advances it, and whether the restriction is no more extensive than necessary (\href{https://supreme.justia.com/cases/federal/us/447/557/}{Justia}).

\section{Detailed Case Facts}

In December 1973, during the Mideastern oil embargo, the New York Public Service Commission (PSC) ordered electric utilities to stop all advertising ``promot[ing] the use of electricity,'' citing insufficient fuel stocks for the winter. After the shortage subsided, the PSC sought comments on keeping the ban. Central Hudson Gas \& Electric Corp., a regulated utility with a geographic monopoly, objected on First Amendment grounds.

In 1977 the PSC extended the ban indefinitely. It prohibited \emph{promotional} advertising ``intended to stimulate the purchase of utility services'' as contrary to national conservation policy, but allowed \emph{institutional and informational} advertising, such as urging customers to shift use to off-peak hours. Denying rehearing, the PSC added that promotional ads would worsen rate inequities, because rates were not based on marginal cost, and would send ``misleading signals'' during an energy crisis.

\section{Procedural History}

Central Hudson sued in New York state court. The trial court, the Appellate Division, and the New York Court of Appeals upheld the order, reasoning that promotional advertising by a monopoly in a ``noncompetitive market'' had little informational value and was outweighed by the State's interest in conservation. The Supreme Court noted probable jurisdiction and reversed on June 20, 1980.

\section{Judicial Votes}

Reversed, \textbf{8--1}.
\begin{itemize}
\item \textbf{Majority:} Justice Powell, joined by Chief Justice Burger and Justices Stewart, White, and Marshall.
\item \textbf{Concurring in the judgment:} Justice Brennan; Justice Blackmun, joined by Justice Brennan; Justice Stevens, joined by Justice Brennan.
\item \textbf{Dissent:} Justice Rehnquist.
\end{itemize}

\section{Holding}

The Constitution gives commercial speech lesser protection than other expression, but restrictions on truthful commercial speech about lawful activity must satisfy a four-part intermediate-scrutiny test. New York's complete ban fails the fourth prong because it is more extensive than necessary to serve the State's substantial interest in energy conservation.

\section{Analysis of Opinions}

\textbf{Majority (Powell).} Synthesizing the cases since \textit{Virginia Pharmacy Board} (1976), Powell held that commercial speech draws its value from its informational function and is hardier than political speech, so false or unlawful messages are unprotected. Prong 1: the ads were truthful and lawful; utilities face interfuel competition, and even a monopolist's ads inform consumers. Prong 2: conservation and fair rates were ``clear and substantial.'' Prong 3: the rate-equity link was ``highly speculative,'' but there was an ``immediate connection'' between advertising and demand. Prong 4: the total ban barred ads for energy-efficient devices such as heat pumps, and the PSC had not shown that narrower measures, such as efficiency disclosures or previewing campaigns, would fail.

\textbf{Blackmun (concurring in the judgment).} The test is too weak. Suppressing truthful information to ``dampen'' demand is covert paternalism that ``strikes at the heart of the First Amendment''; the State should regulate consumption directly rather than ``keeping the public in ignorance.''

\textbf{Stevens (concurring in the judgment).} Powell's definition of commercial speech as ``expression related solely to the economic interests of the speaker and its audience'' is dangerously broad, and the ban suppressed public-issue advocacy.

\textbf{Brennan} joined both Stevens and Blackmun.

\textbf{Rehnquist (dissenting).} A state-created monopoly should not claim immunity from conservation orders; elevating advertising resurrects \textit{Lochner}; and the fourth prong invites judges to second-guess utility engineering.

\textbf{Later developments.} \textit{SUNY v. Fox} (1989) required only a ``reasonable fit,'' \textit{44 Liquormart} (1996) echoed Blackmun's anti-paternalism, and \textit{Sorrell v. IMS Health} (2011) applied ``heightened judicial scrutiny'' to content- and speaker-based rules.

\section{Examples: Future Applications}

\begin{enumerate}
\item \textbf{Same-Side (Invalid): E-Cigarette Ads.} A total ban on truthful adult-oriented e-cigarette ads likely fails the fit prong under \textit{Lorillard} and \textit{44 Liquormart}, given age-gating, warnings, or direct regulation.
\item \textbf{Same-Side (Invalid): Liquor Prices.} A ban on advertising liquor prices seeks to influence conduct by keeping consumers uninformed; \textit{44 Liquormart} struck down a similar ban.
\item \textbf{Opposite-Side (Valid): Illegal Betting Ads.} Ads urging residents to place bets that are crimes in their State concern unlawful activity and fail the first prong.
\item \textbf{Opposite-Side (Valid): Lawyer Solicitation.} A 30-day bar on targeted letters to accident victims, supported by survey evidence, was upheld in \textit{Florida Bar v. Went For It} (1995).
\item \textbf{Fence-Sitter: Marijuana Billboards.} A legalizing State bans dispensary billboards; whether ``lawful activity'' turns on state or federal law is unresolved.
\end{enumerate}

\section{Critique}

\textbf{The boundary problem.} Powell offered two inconsistent definitions of commercial speech: ``expression related solely to the economic interests of the speaker and its audience'' and ``speech proposing a commercial transaction.'' \textit{Bolger} (1983) treats the second as the core, but mixed political-commercial corporate speech still blurs the line.

\textbf{First Amendment Lochnerism.} Rehnquist's dissent presaged critiques that aggressive commercial-speech review, as in \textit{Sorrell} and \textit{44 Liquormart}, lets courts deregulate consumer protection, data privacy, and public health in the name of free speech.

\section{Key Quotations}

\begin{quote}
``In commercial speech cases, then, a four-part analysis has developed. At the outset, we must determine whether the expression is protected by the First Amendment. For commercial speech to come within that provision, it at least must concern lawful activity and not be misleading. Next, we ask whether the asserted governmental interest is substantial. If both inquiries yield positive answers, we must determine whether the regulation directly advances the governmental interest asserted, and whether it is not more extensive than is necessary to serve that interest.''
\end{quote}
\noindent --- Justice Powell, for the Court, 447 U.S. at 566.

\begin{quote}
``The Constitution therefore accords a lesser protection to commercial speech than to other constitutionally guaranteed expression.''
\end{quote}
\noindent --- Justice Powell, for the Court, 447 U.S. at 562--563.

\begin{quote}
``No differences between commercial speech and other protected speech justify suppression of commercial speech in order to influence public conduct through manipulation of the availability of information.''
\end{quote}
\noindent --- Justice Blackmun, concurring in the judgment, 447 U.S. at 578.

\end{document}
